\section{Lecture 4: Complex Predictors, Transformations, and Sampling Distributions}

\subsection{Interaction}

An interaction occurs when a predictor has a different effect on the response depending on the value of another predictor. The usual way of including interactions in a linear model is to add a product term:
\[
Y = \beta_0 + \beta_1 x_1 + \beta_2 x_2 + \beta_3 x_1 x_2 + \varepsilon.
\]

The coefficient $\beta_3$ is estimated in exactly the same way as any other regression coefficient, using ordinary least squares.

\subsection{Diagnostics and Terminology}

Diagnostics for a product term are the same as for any other predictor. The residuals should have the same distribution regardless of the values of $x_i x_j$, and standard diagnostic plots can be made using $x_i x_j$ as the predictor.

The coefficients $\beta_1$ and $\beta_2$ are called the \emph{main effects}, while $\beta_3$ is called the \emph{interaction effect}.

The fitted model is
\[
\widehat{\E}(Y \mid X = x) = \hat{Y}
= \bhat{\beta}_0 + \bhat{\beta}_1 x_1 + \bhat{\beta}_2 x_2 + \bhat{\beta}_3 x_1 x_2.
\]

\subsection{Interpretation of a Continuous--Continuous Interaction}

Fix all other variables. The effect of one variable depends on the value of another fixed variable.

Given $x_2 = a$:
\[
\hat{Y}
= \bhat{\beta}_0 + \bhat{\beta}_2 a + (\bhat{\beta}_1 + a \bhat{\beta}_3) x_1.
\]

Here, $\bhat{\beta}_0 + \bhat{\beta}_2 a$ is the intercept, and $\bhat{\beta}_1 + a \bhat{\beta}_3$ is the slope with respect to $x_1$.

\textbf{Interpretation:} Given $x_2 = a$, the average change in the response for a one-unit increase in $x_1$, holding all predictors fixed, is
\[
\bhat{\beta}_1 + a \bhat{\beta}_3.
\]

\subsection{Interactions Between Indicators and a Continuous Predictor}

Each continuous predictor has a different slope for each level of a categorical variable. Consider
\[
\hat{Y}
= \bhat{\beta}_0
+ \bhat{\beta}_1 I(\text{category B})
+ \bhat{\beta}_2 I(\text{category C})
+ \bhat{\beta}_3 x_3
+ \bhat{\beta}_4 I(\text{category B})\,x_3
+ \bhat{\beta}_5 I(\text{category C})\,x_3.
\]

Each category has a different intercept, and each category has a different slope for $x_3$.

Fixing all other variables, and given category B:
\[
\hat{Y}
= \bhat{\beta}_0 + \bhat{\beta}_1 + (\bhat{\beta}_3 + \bhat{\beta}_4) x_3.
\]

\textbf{Interpretation:} For observations in category B, the intercept is $\bhat{\beta}_0 + \bhat{\beta}_1$, and the average change in the response for a one-unit increase in $x_3$ is $\bhat{\beta}_3 + \bhat{\beta}_4$. This confirms that the effect of $x_3$ depends on category membership, which is the defining feature of an interaction.

Fixing all other variables, and given category A (the reference category):
\[
\hat{Y}
= \bhat{\beta}_0 + \bhat{\beta}_3 x_3.
\]

\textbf{Interpretation:} For observations in category A, the intercept is $\bhat{\beta}_0$, and the average change in the response for a one-unit increase in $x_3$ is $\bhat{\beta}_3$.

\subsection{Polynomial Regression}

We can add new predictors as polynomial powers of an existing predictor variable. A polynomial regression model of degree $k$ is
\[
Y = \beta_0 + \beta_1 x + \beta_2 x^2 + \cdots + \beta_k x^k + \varepsilon.
\]

This model is still linear in the parameters when we define
\[
x_1 = x, \quad x_2 = x^2, \quad \ldots, \quad x_k = x^k.
\]

However, the usual interpretation of slope coefficients as constant marginal effects no longer applies when polynomial terms are included.

\subsection{Transformations}

We usually consider transformations when diagnostic plots, such as residual plots, provide strong evidence against the adequacy of the regression model.

The main reasons for applying transformations are to overcome:
\begin{itemize}
  \item Problems due to non-linearity.
  \item Problems due to non-constant variance.
\end{itemize}

\subsection{General Transformation Framework}

If $Y$ and $X$ are the observed variables, we may fit the model
\[
g_y(Y) = \beta_0 + \beta_1 g_x(X) + \varepsilon,
\]
where $g_y(\cdot)$ and $g_x(\cdot)$ are transformations applied to the response and predictor, respectively.

\subsection{Common Transformations}

Commonly used transformations include:
\[
\text{Logarithmic: } g(x) = \log(x),
\]
\[
\text{Power: } g(x) = x^\lambda \quad (\lambda > 0),
\]
\[
\text{Inverse power: } g(x) = \frac{1}{x^\lambda} \quad (\lambda > 0).
\]

The function $g$ may represent a transformation of the predictor $g_x$ or of the response $g_y$.

\subsubsection{Transforming the Predictors Only}

If the response is not transformed, then $g_y(Y) = Y$, and the model becomes
\[
Y = \beta_0 + \beta_1 g_x(X) + \varepsilon.
\]

Examples:
\[
g_x(x) = \log(x) \quad \Rightarrow \quad Y = \beta_0 + \beta_1 \log(x) + \varepsilon,
\]
\[
g_x(x) = x^\lambda \quad \Rightarrow \quad Y = \beta_0 + \beta_1 x^\lambda + \varepsilon,
\]
\[
g_x(x) = \exp(x) \quad \Rightarrow \quad Y = \beta_0 + \beta_1 \exp(x) + \varepsilon.
\]

\subsubsection{Transforming the Response Only}

If the predictor is not transformed, then $g_x(X) = X$, and the model becomes
\[
g_y(Y) = \beta_0 + \beta_1 X + \varepsilon.
\]

Examples:
\[
g_y(y) = \log(y) \quad \Rightarrow \quad \log(Y) = \beta_0 + \beta_1 X + \varepsilon,
\]
\[
g_y(y) = y^\lambda \quad \Rightarrow \quad Y^\lambda = \beta_0 + \beta_1 X + \varepsilon,
\]
\[
g_y(y) = \exp(y) \quad \Rightarrow \quad \exp(Y) = \beta_0 + \beta_1 X + \varepsilon.
\]

\subsubsection{Transforming Both Predictors and Response}

In the most general case, both the response and predictor are transformed:
\[
g_y(Y) = \beta_0 + \beta_1 g_x(X) + \varepsilon.
\]

Examples:
\[
g(\cdot) = \log(\cdot) \quad \Rightarrow \quad \log(Y) = \beta_0 + \beta_1 \log(X) + \varepsilon,
\]
\[
g(\cdot) = (\cdot)^\lambda \quad \Rightarrow \quad Y^\lambda = \beta_0 + \beta_1 X^\lambda + \varepsilon,
\]
\[
g(\cdot) = \exp(\cdot) \quad \Rightarrow \quad \exp(Y) = \beta_0 + \beta_1 \exp(X) + \varepsilon.
\]

More specifically, we may use different transformations for $Y$ and $X$:
\[
g_y(Y) = \beta_0 + \beta_1 g_x(x) + \varepsilon.
\]

Examples:
\begin{itemize}
  \item $g_y(y) = \log(y)$ and $g_x(x) = \sqrt{x}$
  $\quad \Rightarrow \quad \log(Y) = \beta_0 + \beta_1 \sqrt{x} + \varepsilon$.

  \item $g_y(y) = y^\lambda$ and $g_x(x) = \log(x)$
  $\quad \Rightarrow \quad Y^\lambda = \beta_0 + \beta_1 \log(x) + \varepsilon$.

  \item $g_y(y) = \exp(y)$ and $g_x(x) = \sin\!\left(\tfrac{3\pi}{4}x\right)$
  $\quad \Rightarrow \quad \exp(Y) = \beta_0 + \beta_1 \sin\!\left(\tfrac{3\pi}{4}x\right) + \varepsilon$.
\end{itemize}

\subsection{Choosing a Transformation}

To choose an appropriate transformation:
\begin{itemize}
  \item Plot the response versus predictors to assess potential transformations.
  \item Plot residuals versus fitted values and versus predictors.
  \item Choose a transformation that improves the residual plots.
  \item When both the response $Y$ and the predictor $x$ are measured in the same units, it is often natural to consider the same transformation for both $x$ and $Y$.
\end{itemize}

\subsection{Advantages and Drawbacks of Transformations}

Transformations may make interpretation of the estimated coefficients more difficult. However, transformations often improve predictive performance.

\subsection{Box--Cox Transformation Method}

The Box--Cox transformation attempts to ``fix'' the residuals by transforming the response variable so that the resulting residuals are approximately normally distributed. The transformation is defined as
\[
\psi(Y,\lambda)=
\begin{cases}
\mathrm{gm}(Y)^{\,1-\lambda}\,\dfrac{Y^{\lambda}-1}{\lambda}, & \text{if } \lambda \neq 0, \\[1em]
\mathrm{gm}(Y)\,\log(Y), & \text{if } \lambda = 0,
\end{cases}
\]
where
\[
\mathrm{gm}(Y)
= \prod_{i=1}^{n} Y_i^{1/n}
= \exp\!\left(\frac{1}{n}\sum_{i=1}^{n}\log(Y_i)\right)
\]
is the geometric mean of $Y$.

The Box--Cox transformation uses maximum likelihood to choose the value of $\lambda$ so that the residuals are approximately normally distributed. It is an automatic, likelihood-based method for selecting a transformation and provides a confidence interval for plausible values of $\lambda$.

The Box--Cox transformation is \emph{not} intended to enforce linearity. Instead, it is used to address non-normality and non-constant variance by transforming the response variable, which indirectly improves the behavior of the residuals.

\subsection{Interpreting Transformations of Response and Predictors}

To interpret regression models that use transformations, always remember to incorporate the variables exactly as they appear in the model, such as $\log(Y)$, $\log(x)$, and similar transformed forms.

If a predictor appears squared in the model, then the interpretation is phrased as ``for a one-unit increase in $x^2$.''

When the response variable has been transformed, predicted values represent conditional means on the transformed scale, not on the original measurement scale.

\textbf{Interpretation on the Original Scale:} In general, we do not back-transform fitted values to obtain predictions on the original measurement scale. See Weisberg (2005) and Sheather (2009) for examples where certain transformations admit approximate interpretations on the original scale.

\includegraphics[width = \textwidth]{pictures/boxcox.png}

\subsection{Review: Expectation and Covariance of a Random Vector}

Let $\mathbf{Y} = (Y_1, \dots, Y_n)^{\top}$ be a random column vector.

The expected value of $\mathbf{Y}$ is defined as
\[
\E(\mathbf{Y})
=
(\E(Y_1), \dots, \E(Y_n))^{\top}.
\]

The covariance matrix (sometimes called the variance--covariance matrix) of $\mathbf{Y}$ is defined as
\[
\Var(\mathbf{Y})
=
\Cov(\mathbf{Y})
=
\E\!\left[(\mathbf{Y}-\E(\mathbf{Y}))(\mathbf{Y}-\E(\mathbf{Y}))^{\top}\right].
\]

The matrix $\Var(\mathbf{Y})$ has dimensions $n \times n$. The $(i,j)$-th entry of $\Var(\mathbf{Y})$ represents the covariance between $Y_i$ and $Y_j$.

\textbf{Example ($n=3$):}
\[
\E(\mathbf{Y})
=
\begin{pmatrix}
\E(Y_1) \\
\E(Y_2) \\
\E(Y_3)
\end{pmatrix},
\qquad
\Var(\mathbf{Y})
=
\begin{pmatrix}
\Var(Y_1) & \Cov(Y_1,Y_2) & \Cov(Y_1,Y_3) \\
\Cov(Y_2,Y_1) & \Var(Y_2) & \Cov(Y_2,Y_3) \\
\Cov(Y_3,Y_1) & \Cov(Y_3,Y_2) & \Var(Y_3)
\end{pmatrix}.
\]

Note that $\Var(Y_1) = \Cov(Y_1,Y_1)$.

For two random variables:
\[
\Cov(Y_1,Y_2) = \E(Y_1Y_2) - \E(Y_1)\E(Y_2).
\]

\begin{prop}{Properties of Expectation and Variance}{prop:evprop}
Let $\mathbf{Y}$ be a random column vector of dimension $n$ with mean $\boldsymbol{\mu}$ and covariance matrix $\boldsymbol{\Sigma}$. Let $\mathbf{A}$ be a real matrix of dimension $n \times p$, and let $\mathbf{a} \in \mathbb{R}^p$. Then:
\[
\Var(\mathbf{Y}) = \E(\mathbf{Y}\mathbf{Y}^{\top}) - \boldsymbol{\mu}\boldsymbol{\mu}^{\top},
\]
\[
\E(\mathbf{A}^{\top}\mathbf{Y} + \mathbf{a}) = \mathbf{A}^{\top}\boldsymbol{\mu} + \mathbf{a},
\]
\[
\Var(\mathbf{A}^{\top}\mathbf{Y} + \mathbf{a}) = \mathbf{A}^{\top}\boldsymbol{\Sigma}\mathbf{A}.
\]
\end{prop}

\clearpage

\section*{Sampling Distributions of the Least Squares Estimators}

Consider the multiple linear regression model:
\[
\mathbf{Y} = \mathbf{X}\bbeta + \beps,
\qquad
\beps \sim \mathcal{N}_n(\mathbf{0}, \sigma^2 \mathbf{I}_n).
\]

Then
\[
\E(\beps)
= 0 =
\begin{pmatrix}
0 \\[0.2em]
0 \\[0.2em]
\vdots \\[0.2em]
0
\end{pmatrix},
\qquad
\Var(\beps)
= \sigma^2 \mathbf{I}_n =
\begin{pmatrix}
\sigma^2 & 0 & \cdots & 0 \\
0 & \sigma^2 & \cdots & 0 \\
\vdots & \vdots & \ddots & \vdots \\
0 & 0 & \cdots & \sigma^2
\end{pmatrix}.
\]

Thus:
\[
\Var(\varepsilon_i) = \sigma^2,
\qquad \text{for } i = 1, \dots, n,
\]
\[
\Cov(\varepsilon_i, \varepsilon_j) = 0,
\qquad \text{for } i \neq j.
\]

It follows that
\[
\E(\mathbf{Y}) = \mathbf{X}\bbeta,
\qquad
\Var(\mathbf{Y}) = \sigma^2 \mathbf{I}_n,
\]
and hence
\[
\mathbf{Y} \sim \mathcal{N}_n(\mathbf{X}\bbeta, \sigma^2 \mathbf{I}_n).
\]

The least squares estimator of $\bbeta$ is
\[
\bhat{\beta} = (\mathbf{X}^\top \mathbf{X})^{-1}\mathbf{X}^\top \mathbf{Y},
\]
which is a random vector.

\subsection{Unbiasedness of the Least Squares Estimator}

Using the model $\mathbf{Y} = \mathbf{X}\bbeta + \beps$:

\[
\bhat{\beta} - \bbeta
= (\mathbf{X}^\top \mathbf{X})^{-1}\mathbf{X}^\top \mathbf{Y} - \bbeta
= (\mathbf{X}^\top \mathbf{X})^{-1}\mathbf{X}^\top (\mathbf{X}\bbeta + \beps) - \bbeta
= (\mathbf{X}^\top \mathbf{X})^{-1}\mathbf{X}^\top \beps.
\]

Conditioning on $\mathbf{X}$:
\[
\E(\bhat{\beta} - \bbeta \mid \mathbf{X})
=
(\mathbf{X}^\top \mathbf{X})^{-1}\mathbf{X}^\top \E(\beps)
=
\mathbf{0},
\]
since $\E(\beps) = \mathbf{0}$.

By linearity of expectation:
\[
\E(\bhat{\beta} \mid \mathbf{X})
=
\E(\bbeta \mid \mathbf{X})
=
\bbeta.
\]

Thus, $\bhat{\beta}$ is an unbiased estimator of $\bbeta$.

\subsection{Variance of the Least Squares Estimator}

Using the property $\Var(\mathbf{A}\mathbf{Y}) = \mathbf{A}\,\Var(\mathbf{Y})\,\mathbf{A}^\top$:
\[
\Var(\bhat{\beta} - \bbeta \mid \mathbf{X})
=
\Var\!\left( (\mathbf{X}^\top \mathbf{X})^{-1}\mathbf{X}^\top \beps \mid \mathbf{X} \right).
\]

Therefore:
\[
\Var(\bhat{\beta} - \bbeta \mid \mathbf{X})
=
(\mathbf{X}^\top \mathbf{X})^{-1}\mathbf{X}^\top
\Var(\beps)
\mathbf{X}(\mathbf{X}^\top \mathbf{X})^{-1}.
\]

Substituting $\Var(\beps) = \sigma^2 \mathbf{I}_n$:
\[
\Var(\bhat{\beta} - \bbeta \mid \mathbf{X})
=
\sigma^2(\mathbf{X}^\top \mathbf{X})^{-1}\mathbf{X}^\top \mathbf{X}(\mathbf{X}^\top \mathbf{X})^{-1}
=
\sigma^2(\mathbf{X}^\top \mathbf{X})^{-1}.
\]

Hence:
\[
\Var(\bhat{\beta} \mid \mathbf{X})
=
\sigma^2(\mathbf{X}^\top \mathbf{X})^{-1}.
\]

\begin{keyresult}[Sampling Distribution of the Least Squares Estimator]
\[
\bhat{\beta}
=
\begin{pmatrix}
\bhat{\beta}_0 \\
\bhat{\beta}_1 \\
\vdots \\
\bhat{\beta}_p
\end{pmatrix}
\sim
\mathcal{N}_{p+1}\!\left(
\bbeta,\,
\sigma^2(\mathbf{X}^\top \mathbf{X})^{-1}
\right).
\]

Under the assumption $\beps \sim \mathcal{N}_n(\mathbf{0}, \sigma^2 \mathbf{I}_n)$:
\[
\beps \sim \mathcal{N}_n\!\left(\mathbf{0},\, \sigma^2 \mathbf{I}_n\right).
\]
\end{keyresult}

\subsection{Simple Linear Regression: Sampling Distribution of \texorpdfstring{$\bhat{\beta}$}{beta-hat}}

Assume the simple linear regression model (with intercept):
\[
Y_i = \beta_0 + \beta_1 x_i + \varepsilon_i,
\qquad i=1,\dots,n,
\]
or in matrix form:
\[
\mathbf{Y}=\mathbf{X}\bbeta+\beps,
\qquad
\beps\sim \mathcal{N}_n(\mathbf{0},\sigma^2\mathbf{I}_n),
\]
where
\[
\mathbf{Y}=
\begin{pmatrix}
Y_1\\ \vdots\\ Y_n
\end{pmatrix},
\quad
\bbeta=
\begin{pmatrix}
\beta_0\\ \beta_1
\end{pmatrix},
\quad
\mathbf{X}=
\begin{pmatrix}
1 & x_1\\
\vdots & \vdots\\
1 & x_n
\end{pmatrix}.
\]

\subsubsection{Step 1: Distribution of \texorpdfstring{$\bhat{\beta}$}{beta-hat}}

The least squares estimator is
\[
\bhat{\beta}
=
(\mathbf{X}^\top \mathbf{X})^{-1}\mathbf{X}^\top \mathbf{Y}.
\]

Substitute $\mathbf{Y}=\mathbf{X}\bbeta+\beps$:
\[
\bhat{\beta}
=
(\mathbf{X}^\top \mathbf{X})^{-1}\mathbf{X}^\top(\mathbf{X}\bbeta+\beps)
=
\bbeta+(\mathbf{X}^\top \mathbf{X})^{-1}\mathbf{X}^\top\beps.
\]

Hence:
\[
\bhat{\beta}-\bbeta
=
(\mathbf{X}^\top \mathbf{X})^{-1}\mathbf{X}^\top\beps.
\]

Since $\beps\sim \mathcal{N}_n(\mathbf{0},\sigma^2\mathbf{I}_n)$ and a linear transformation of a multivariate normal is multivariate normal:
\[
\bhat{\beta} \sim \mathcal{N}_2\!\left(\bbeta,\ \sigma^2(\mathbf{X}^\top\mathbf{X})^{-1}\right).
\]

\subsubsection{Step 2: Compute $\mathbf{X}^\top\mathbf{X}$ and its Inverse}

Compute $\mathbf{X}^\top$:
\[
\mathbf{X}^\top=
\begin{pmatrix}
1 & \cdots & 1\\
x_1 & \cdots & x_n
\end{pmatrix}.
\]

Therefore:
\[
\mathbf{X}^\top\mathbf{X}
=
\begin{pmatrix}
\sum_{i=1}^n 1 & \sum_{i=1}^n x_i\\[0.4em]
\sum_{i=1}^n x_i & \sum_{i=1}^n x_i^2
\end{pmatrix}
=
\begin{pmatrix}
n & \sum_{i=1}^n x_i\\[0.4em]
\sum_{i=1}^n x_i & \sum_{i=1}^n x_i^2
\end{pmatrix}.
\]

Its determinant is:
\[
\det(\mathbf{X}^\top\mathbf{X})
=
n\sum_{i=1}^n x_i^2-\left(\sum_{i=1}^n x_i\right)^2.
\]

Using $\bar{x}=\frac{1}{n}\sum_{i=1}^n x_i$:
\[
\det(\mathbf{X}^\top\mathbf{X})
=
n\sum_{i=1}^n x_i^2-n^2\bar{x}^2
=
n\left(\sum_{i=1}^n x_i^2-n\bar{x}^2\right).
\]

Define
\[
s_{xx}=\sum_{i=1}^n (x_i-\bar{x})^2
=\sum_{i=1}^n x_i^2-2\bar{x}\sum_{i=1}^n x_i+n\bar{x}^2
=\sum_{i=1}^n x_i^2-n\bar{x}^2,
\]
so
\[
\det(\mathbf{X}^\top\mathbf{X})=n\,s_{xx}.
\]

For a $2\times 2$ matrix:
\[
\begin{pmatrix}
a & b\\
b & d
\end{pmatrix}^{-1}
=
\frac{1}{ad-b^2}
\begin{pmatrix}
d & -b\\
-b & a
\end{pmatrix}.
\]

Applying this with $a=n$, $b=\sum x_i$, $d=\sum x_i^2$:
\[
(\mathbf{X}^\top\mathbf{X})^{-1}
=
\frac{1}{n\sum_{i=1}^n x_i^2-(\sum_{i=1}^n x_i)^2}
\begin{pmatrix}
\sum_{i=1}^n x_i^2 & -\sum_{i=1}^n x_i\\[0.4em]
-\sum_{i=1}^n x_i & n
\end{pmatrix}.
\]

Rewrite in terms of $\bar{x}$ and $s_{xx}$. Since $\sum x_i=n\bar{x}$ and $\sum x_i^2 = s_{xx}+n\bar{x}^2$:
\[
(\mathbf{X}^\top\mathbf{X})^{-1}
=
\frac{1}{n s_{xx}}
\begin{pmatrix}
s_{xx}+n\bar{x}^2 & -n\bar{x}\\[0.4em]
-n\bar{x} & n
\end{pmatrix}
=
\frac{1}{s_{xx}}
\begin{pmatrix}
\frac{1}{n}\sum_{i=1}^n x_i^2 & -\bar{x}\\[0.4em]
-\bar{x} & 1
\end{pmatrix}.
\]

\subsubsection{Step 3: Variance Matrix and Individual Variances}

From Step 1:
\[
\Var(\bhat{\beta})
=
\sigma^2(\mathbf{X}^\top\mathbf{X})^{-1}
=
\frac{\sigma^2}{s_{xx}}
\begin{pmatrix}
\frac{1}{n}\sum_{i=1}^n x_i^2 & -\bar{x}\\[0.4em]
-\bar{x} & 1
\end{pmatrix}.
\]

Thus:
\[
\Var(\bhat{\beta}_0)
=
\frac{\sigma^2}{s_{xx}}\left(\frac{1}{n}\sum_{i=1}^n x_i^2\right),
\qquad
\Var(\bhat{\beta}_1)
=
\frac{\sigma^2}{s_{xx}},
\qquad
\Cov(\bhat{\beta}_0,\bhat{\beta}_1)
=
-\frac{\sigma^2\bar{x}}{s_{xx}}.
\]

\begin{keyresult}[Simple Linear Regression Sampling Distribution]
Since $p=1$ in simple linear regression (intercept + one slope):
\[
\bhat{\beta}
=
\begin{pmatrix}
\bhat{\beta}_0\\
\bhat{\beta}_1
\end{pmatrix}
\sim
\mathcal{N}_2\!\left(
\begin{pmatrix}
\beta_0\\
\beta_1
\end{pmatrix},
\ \frac{\sigma^2}{s_{xx}}
\begin{pmatrix}
\frac{1}{n}\sum_{i=1}^n x_i^2 & -\bar{x}\\[0.4em]
-\bar{x} & 1
\end{pmatrix}
\right),
\]
where $s_{xx}=\sum_{i=1}^n (x_i-\bar{x})^2$.
\end{keyresult}

\subsection{Estimated Standard Errors}

\[
\se(\bhat{\beta}_0) = \sqrt{\Var(\bhat{\beta}_0)}.
\]

\[
\widehat{\se}(\bhat{\beta}_0)
= \hat{\sigma}\,
\sqrt{
\frac{1}{n}
+
\frac{\bar{x}^2}{s_{xx}}
}.
\]

\[
\se(\bhat{\beta}_1) = \sqrt{\Var(\bhat{\beta}_1)}.
\]

\[
\widehat{\se}(\bhat{\beta}_1)
=  \hat{\sigma}\,
\sqrt{
\frac{1}{s_{xx}}
}.
\]

where
\[
\hat{\sigma}^2
=
\frac{1}{\,n-(p+1)\,}
\sum_{i=1}^n \left( Y_i - \hat{Y}_i \right)^2 = \frac{\RSS}{n-2}.
\]

\subsection{Simple Linear Regression for the Iris Data}

We consider the simple linear regression model:
\[
\text{Petal}_L = \beta_0 + \beta_1 \,\text{Sepal}_L + \varepsilon,
\qquad \varepsilon \sim \mathcal{N}(0,\sigma^2).
\]

The following \texttt{R} code is used to compute the fitted model quantities and the estimated standard errors:
\begin{verbatim}
data(iris)
fit <- lm(Petal.Length ~ Sepal.Length, data = iris)
fitted_values <- fitted(fit)

y <- iris$Petal.Length
x <- iris$Sepal.Length
mean_x <- mean(x)

RSS <- sum((y - fitted_values)^2)
n.iris <- nrow(iris)

s_xx <- sum((x - mean_x)^2)
\end{verbatim}

The model summary output is:
\begin{verbatim}
Coefficients:
              Estimate Std. Error t value Pr(>|t|)
(Intercept)   -7.10144   0.50666   -14.02  <2e-16
Sepal.Length   1.85843   0.08586    21.65  <2e-16
\end{verbatim}

The estimated standard errors of the least squares estimators are:
\[
\se(\bhat{\beta}_0) = 0.50666,
\qquad
\se(\bhat{\beta}_1) = 0.08586.
\]

\subsection{Multiple Linear Regression for the Iris Data}

We now consider the multiple linear regression model:
\[
\text{Petal}_L
= \beta_0 + \beta_1 \,\text{Sepal}_L + \beta_2 \,\text{Sepal}_W + \varepsilon,
\qquad \varepsilon \sim \mathcal{N}(0,\sigma^2).
\]

The corresponding \texttt{R} code is:
\begin{verbatim}
data(iris)
fit <- lm(Petal.Length ~ Sepal.Length + Sepal.Width, data = iris)
summary(fit)
\end{verbatim}

The coefficient table from the model output is:
\begin{verbatim}
Coefficients:
              Estimate Std. Error t value Pr(>|t|)
(Intercept)   -2.52476   0.56344   -4.481  1.48e-05
Sepal.Length   1.77559   0.06441    27.569  <2e-16
Sepal.Width   -1.33862   0.12236  -10.940  <2e-16
\end{verbatim}

The estimated standard errors of the least squares estimators are:
\[
\se(\bhat{\beta}_0) = 0.56344,
\qquad
\se(\bhat{\beta}_1) = 0.06441,
\qquad
\se(\bhat{\beta}_2) = 0.12236.
\]

\subsection{Interpretation}

Including $\text{Sepal}_W$ in the model reduces the standard error of $\bhat{\beta}_1$ (from $0.08586$ to $0.06441$), indicating increased precision in estimating the effect of sepal length once sepal width is accounted for.

\pagebreak

\section*{Tutorial 1 (January 28)}

\subsection{Problem 1}

Sometimes it is convenient to write the simple linear regression model in a different form. Starting from
\[
Y_i = \beta_0 + \beta_1 x_i + \varepsilon_i,
\]
add and subtract $\beta_1 \bar{x}$ to obtain
\[
Y_i = \beta_0 + \beta_1 \bar{x} + \beta_1(x_i - \bar{x}) + \varepsilon_i
= \alpha + \beta_1(x_i - \bar{x}) + \varepsilon_i,
\]
where $\alpha = \beta_0 + \beta_1 \bar{x}$. This is called the \textbf{deviations-from-the-mean form}.

\subsubsection*{Solution}

\textbf{(1) Meaning of $\alpha$:}

The parameter $\alpha$ represents the expected value of $Y_i$ when $x_i = \bar{x}$. That is:
\[
\E[Y_i \mid x_i = \bar{x}] = \alpha.
\]

\textbf{(2) Least squares estimates:}

The model is
\[
Y_i = \alpha + \beta_1(x_i - \bar{x}) + \varepsilon_i.
\]

The normal equations give
\[
\sum_{i=1}^n (Y_i - \hat{\alpha} - \hat{\beta}_1(x_i - \bar{x})) = 0,
\]
which implies
\[
\hat{\alpha} = \bar{y}.
\]

For $\beta_1$:
\[
\sum_{i=1}^n (x_i - \bar{x})(Y_i - \hat{\alpha} - \hat{\beta}_1(x_i - \bar{x})) = 0,
\]
leading to
\[
\bhat{\beta}_1 =
\frac{\sum_{i=1}^n (x_i - \bar{x})(y_i - \bar{y})}
{\sum_{i=1}^n (x_i - \bar{x})^2}.
\]

\subsection{Problem 2}

Consider the simple linear regression model:
\[
Y_i = \beta x_i + \varepsilon_i, \quad i = 1,\dots,n.
\]

Show that the least squares estimate of $\beta$ is
\[
\bhat{\beta} = \frac{\sum_{i=1}^n x_i y_i}{\sum_{i=1}^n x_i^2}.
\]

\subsubsection*{Solution}

The residual sum of squares is:
\[
S(\beta) = \sum_{i=1}^n (y_i - \beta x_i)^2.
\]

Differentiate with respect to $\beta$ and set to zero:
\[
\frac{dS}{d\beta} = -2\sum_{i=1}^n x_i(y_i - \beta x_i) = 0.
\]

Solving gives:
\[
\bhat{\beta} = \frac{\sum_{i=1}^n x_i y_i}{\sum_{i=1}^n x_i^2}.
\]

\subsection{Problem 3}

Show that
\[
(y - \mathbf{X}\bhat{\beta})^\top (y - \mathbf{X}\bhat{\beta}) = y^\top y - \bhat{\beta}^\top \mathbf{X}^\top y.
\]

\subsubsection*{Solution}

Expand the left-hand side:
\[
(\bY - \mathbf{X}\bhat{\beta})^\top (\bY - \mathbf{X}\bhat{\beta})
= \bY^\top \bY - 2\bhat{\beta}^\top \mathbf{X}^\top \bY + \bhat{\beta}^\top \mathbf{X}^\top \mathbf{X} \bhat{\beta}.
\]

Using the normal equations $\mathbf{X}^\top \mathbf{X} \bhat{\beta} = \mathbf{X}^\top \bY$:
\[
\bhat{\beta}^\top \mathbf{X}^\top \mathbf{X} \bhat{\beta} = \bhat{\beta}^\top \mathbf{X}^\top \bY.
\]

Substituting yields:
\[
\bY^\top \bY - \bhat{\beta}^\top \mathbf{X}^\top \bY.
\]

\subsection{Problem 4}

Consider the simple linear regression:
\[
Y_i = \beta_0 + \beta_1 x_i + \varepsilon_i.
\]

Show that
\[
\frac{1}{n}\sum_{i=1}^n \hat{y}_i = \frac{1}{n}\sum_{i=1}^n y_i.
\]

\subsubsection*{Solution}

The fitted values are $\hat{y}_i = \bhat{\beta}_0 + \bhat{\beta}_1 x_i$. Thus:
\[
\frac{1}{n}\sum_{i=1}^n \hat{y}_i
= \bhat{\beta}_0 + \bhat{\beta}_1 \bar{x}.
\]

Using $\bhat{\beta}_0 = \bar{y} - \bhat{\beta}_1 \bar{x}$:
\[
\frac{1}{n}\sum_{i=1}^n \hat{y}_i = \bar{y}.
\]

\subsection{Problem 5}

Consider the multiple linear regression:
\[
Y_i = \beta_0 + \beta_1 x_{i1} + \beta_2 x_{i2} + \varepsilon_i.
\]

\subsubsection*{Solution}

\textbf{(1)} The design matrix $\mathbf{X}$ has columns $(1, x_{i1}, x_{i2})$. Therefore:
\[
\mathbf{X}^\top \mathbf{Y} =
\begin{pmatrix}
\sum_{i=1}^n y_i \\
\sum_{i=1}^n x_{i1} y_i \\
\sum_{i=1}^n x_{i2} y_i
\end{pmatrix}.
\]

\textbf{(2)} Similarly:
\[
\mathbf{X}^\top \mathbf{X} =
\begin{pmatrix}
n & \sum_{i=1}^n x_{i1} & \sum_{i=1}^n x_{i2} \\
\sum_{i=1}^n x_{i1} & \sum_{i=1}^n x_{i1}^2 & \sum_{i=1}^n x_{i1}x_{i2} \\
\sum_{i=1}^n x_{i2} & \sum_{i=1}^n x_{i2}x_{i1} & \sum_{i=1}^n x_{i2}^2
\end{pmatrix}.
\]

\clearpage
